\chapter{Supporting Calculations}
\label{app:derivation-details}

These calculations support the variance sensitivity, near-independence
interpretation and validity conditions in Chapter~\ref{chap:expanded}.

\section{Mean of the variance influence function}

The influence function $g_P(x,y)$ is centred at the population level. This
can be shown without expanding its individual terms. The probability-weighted
average of the one-cell perturbation directions is
\begin{equation}
  \sum_{x,y}p_{xy}\{\delta_{(x,y)}-P\}=P-P=0.
\end{equation}
The first derivative is linear in the direction of change. Therefore, the
weighted average of the one-cell derivatives is the derivative in the
weighted-average direction:
\begin{align}
  \E_P\{g_P(X,Y)\}
  &=\sum_{x,y}p_{xy}
    \left.\frac{\mathrm d}{\mathrm d\varepsilon}
    V\{(1-\varepsilon)P+\varepsilon\delta_{(x,y)}\}
    \right|_{0}\\
  &=\left.\frac{\mathrm d}{\mathrm d\varepsilon}
    V\left\{P+\varepsilon
    \sum_{x,y}p_{xy}\{\delta_{(x,y)}-P\}\right\}
    \right|_{0}\\
  &=\left.\frac{\mathrm d}{\mathrm d\varepsilon}V(P)
    \right|_{0}=0.
\end{align}
The same weighted cancellation holds after exact plug-in substitution.

\section{Fixed-margin behaviour near independence}
\label{sec:fixed-margin-local-derivation}

Let \(P_0=(q_{ij})\) be independent, with positive cells
\(q_{ij}=a_i b_j\). Consider \(P_t=P_0+tH\), where every row and column
sum of \(H\) is zero. For sufficiently small \(t\), all cells remain
positive and the margins remain \(a\) and \(b\). We expand around
\(t=0\), where \(I(P_0)=0\).

For \(f(t)=I(P_t)\), the second-order Taylor formula is
\begin{equation}
 f(t)=f(0)+f'(0)t+\frac{f''(0)}{2!}t^2+O(t^3).
\end{equation}
The fixed margins make the pointwise MI in cell \((i,j)\)
\begin{equation}
 \pmi_{P_t}(i,j)
 =\log\left(1+\frac{tH_{ij}}{q_{ij}}\right)
 =\frac{tH_{ij}}{q_{ij}}-
  \frac{t^2H_{ij}^2}{2q_{ij}^2}+O(t^3),
\end{equation}
using \(\log(1+z)=z-z^2/2+O(z^3)\). Multiplying by the cell
probability gives
\begin{equation}
 (q_{ij}+tH_{ij})\pmi_{P_t}(i,j)
 =tH_{ij}+\frac{t^2H_{ij}^2}{2q_{ij}}+O(t^3).
\end{equation}
The first-order terms sum to zero because \(\sum_{i,j}H_{ij}=0\).
Thus \(f(0)=f'(0)=0\), \(f''(0)=\sum_{i,j}H_{ij}^2/q_{ij}\), and
\begin{equation}
 I(P_t)=\frac{t^2}{2}\sum_{i,j}\frac{H_{ij}^2}{q_{ij}}+O(t^3).
\end{equation}

To calculate \(V(P_t)\), square the leading pointwise-MI term. Its
probability-weighted average is
\(t^2\sum_{i,j}H_{ij}^2/q_{ij}+O(t^3)\), while \(I(P_t)^2\) is of order
\(t^4\). Hence
\begin{equation}
 V(P_t)=t^2\sum_{i,j}\frac{H_{ij}^2}{q_{ij}}+O(t^3).
\end{equation}

For the sensitivity in Equation~\eqref{eq:g-final}, the row conditional
mean of pointwise MI has first-order term
\(t\sum_j H_{ij}/a_i=0\); the column mean likewise has first-order term
\(t\sum_i H_{ij}/b_j=0\). The squared centred pointwise MI and \(V(P_t)\)
are both of order \(t^2\). Therefore the bracket in Equation
\eqref{eq:g-final} supplies the first-order term:
\begin{equation}
 g_{P_t}(i,j)=\frac{2tH_{ij}}{q_{ij}}+O(t^2).
\end{equation}
Since \(g_{P_t}\) has mean zero, its probability-weighted variance is
\begin{equation}
 \tau^2(P_t)=4t^2\sum_{i,j}\frac{H_{ij}^2}{q_{ij}}+O(t^3).
\end{equation}
Comparing these three leading terms gives \(V(P_t)\approx2I(P_t)\) and
\(\tau^2(P_t)\approx4V(P_t)\) for small positive MI along this path.
At \(t=0\), the first-order sampling approximation itself degenerates;
Section~\ref{sec:independence-boundary} treats that case separately.

\section{Dependence between the MI and variance estimates}
\label{sec:mi-variance-covariance}

Both \(\Ihat(P)\) and \(\Vhat(P)\) are calculated from the same sample. Under
fixed positive support, their first-order covariance is
\begin{equation}
  \operatorname{Cov}\{\Ihat(P),\Vhat(P)\}
  \approx \frac{1}{n_P}\E_P\!
  \left[\{\pmi_P(X,Y)-I(P)\}g_P(X,Y)\right].
\end{equation}
The corresponding expression for \(Q\) replaces \(P\) and \(n_P\) with
\(Q\) and \(n_Q\). To see its direction near independence, use the fixed-margin
path \(P_t=(q_{ij}+tH_{ij})\) from
Section~\ref{sec:fixed-margin-local-derivation}. For small positive \(t\),
the expansions there give
\begin{equation}
  \pmi_{P_t}(i,j)-I(P_t)
  =\frac{tH_{ij}}{q_{ij}}+O(t^2),
  \qquad
  g_{P_t}(i,j)=\frac{2tH_{ij}}{q_{ij}}+O(t^2).
\end{equation}
Substituting both expressions into the first-order covariance above gives
\begin{align}
  \operatorname{Cov}\{\Ihat(P_t),\Vhat(P_t)\}
  &\approx \frac{2t^2}{n_P}
      \sum_{i,j}\frac{H_{ij}^2}{q_{ij}}
      +O\!\left(\frac{t^3}{n_P}\right)\\
  &=\frac{2V(P_t)}{n_P}
      +O\!\left(\frac{t^3}{n_P}\right).
\end{align}
The first-order covariance is therefore positive for small positive MI
along this path. At exact independence, \(V(P_0)=0\) and a second-order
sampling approximation is needed. This component covariance alone does not
determine the two-sample test's rejection rate.

\section{The independence boundary}
\label{sec:independence-boundary}

At independence, \(p_{ij}=p_{i+}p_{+j}\), so every population
pointwise-MI value is \(\pmi_P(i,j)=\log 1=0\). Consequently,
\(I(P)=V(P)=0\), and the first-order change in MI is zero. Sampling can
still produce positive estimated MI, but its leading variation comes from
second-order cell errors.

To calculate this variation, fix a strictly positive independent table
\(P\) and let \(h=(h_{ij})\) describe a redistribution of probability,
with \(\sum_{i,j}h_{ij}=0\). Define
\begin{equation}
 P_\varepsilon=P+\varepsilon h,
 \qquad f(\varepsilon)=I(P_\varepsilon),
\end{equation}
for \(\varepsilon\) small enough that every cell remains positive.
The expansion point is \(\varepsilon=0\), where \(f(0)=I(P)=0\).
The second-order Taylor formula gives
\begin{equation}
 f(\varepsilon)
 =f(0)+f'(0)\varepsilon+\frac{f''(0)}{2!}\varepsilon^2
  +o(\varepsilon^2).
\end{equation}
The first derivative is
\begin{equation}
 f'(0)=\sum_{i,j}\{\pmi_P(i,j)-1\}h_{ij}
       =-\sum_{i,j}h_{ij}=0.
\end{equation}
Here the pointwise-MI terms vanish at independence, and the remaining
terms cancel because the probabilities must still sum to one.

The row and column changes are \(h_{i+}=\sum_jh_{ij}\) and
\(h_{+j}=\sum_ih_{ij}\). Applying
\(\mathrm d^2(u\log u)/\mathrm du^2=1/u\) to the joint and marginal
terms in Equation~\eqref{eq:mi-probability-sums} gives
\begin{equation}
 f''(0)=\sum_{i,j}\frac{h_{ij}^2}{p_{ij}}
        -\sum_i\frac{h_{i+}^2}{p_{i+}}
        -\sum_j\frac{h_{+j}^2}{p_{+j}}.
\end{equation}
Thus the leading change is
\begin{equation}
 I(P_\varepsilon)
 =\frac{\varepsilon^2}{2}
  \left\{\sum_{i,j}\frac{h_{ij}^2}{p_{ij}}
        -\sum_i\frac{h_{i+}^2}{p_{i+}}
        -\sum_j\frac{h_{+j}^2}{p_{+j}}\right\}
  +o(\varepsilon^2).
 \label{eq:appendix-mi-second-order}
\end{equation}
The first sum measures squared cell changes. The last two remove the
contributions due to changed margins, leaving the change in association.
The remainder is negligible relative to \(\varepsilon^2\) as
\(\varepsilon\) approaches zero.

For a multinomial sample, take \(\varepsilon h=\widehat P-P\).
Cell-probability errors typically have size proportional to \(n^{-1/2}\).
Squaring them makes \(\Ihat(P)\) proportional to \(n^{-1}\) under
independence. Under fixed table dimensions and positive cell probabilities,
the likelihood-ratio result is
\begin{equation}
 2n\Ihat(P)\xrightarrow{d}\chi^2_{(r-1)(c-1)}
 \label{eq:independence-chi-squared}
\end{equation}
\citep{wilks1938,brillinger2004}. This is the usual one-table likelihood-ratio,
or \(G\), test of independence. Its chi-squared distribution comes from
quadratic cell errors, whereas Normal Wald and Expanded Welch use the
first-order MI error. The two-sample first-order approximation therefore
does not apply unchanged at exact independence.

\section{The combined formula when a component is zero}
\label{sec:zero-component}

For positive component degrees of freedom, substituting
Equation~\eqref{eq:nu-p-hat} and its \(Q\) analogue into
Equation~\eqref{eq:expanded-df} gives the equivalent combined form
\begin{equation}
  \widehat\nu_{\mathrm{expanded}}
  =
  \frac{
  \left\{\Vhat(P)/n_P+\Vhat(Q)/n_Q\right\}^2
  }{
  \widehat\tau^2(P)/(2n_P^3)
  +
  \widehat\tau^2(Q)/(2n_Q^3)
  }.
  \label{eq:expanded-df-direct}
\end{equation}
This direct form remains defined when one table has
\(\Vhat=\widehat\tau^2=0\) and the other gives a positive denominator.
The zero contribution drops out, leaving the other table's component degrees
of freedom. The implementation nevertheless requires both components to have
positive, finite degrees of freedom, so it records this case as invalid.

\medskip
\noindent\textbf{Numerical derivative checks.}

The analytic derivatives in Equations~\eqref{eq:mi-derivative} and
\eqref{eq:g-final} were also checked by centred finite differences along
the one-cell perturbation paths in Chapter~\ref{chap:expanded}. Small
positive and negative steps were chosen to keep every cell positive; the
change in MI or variance divided by the distance between the two steps
agreed with its analytic derivative.
